\documentclass{exam}

\newif\ifA
\newif\ifkey
\newif\ifprintselected   % required by exam_config's \ansbox

%%%%%%%%%%%%%%%%%%%%%%%%%%%%%%%%%%%%%%%%%%%%%%%%%%%%%%
%%% SET THESE IF VARIABLES
%%%%%%%%%%%%%%%%%%%%%%%%%%%%%%%%%%%%%%%%%%%%%%%%%%%%%%
% Lab Num -- Module 3's six lab periods follow L04 (Lab 1) and the two
% Module 2 measurement labs, so L19 is Lab 5. Override with \def\labNumIn{}
% if the course sequence renumbers.
\ifdefined\labNumIn
	\def\labNum{\labNumIn}
\else
	\def\labNum{5}
\fi

% is this the key or not?
\ifdefined\iskey
	\ifnum\iskey=1
		\keytrue
	\else
		\keyfalse
	\fi
\else
	%% Pick one manually
	\keyfalse % if this is NOT the key
%	\keytrue     % if this IS the key
\fi

% set up the header and footer - change for every term
\ifdefined\thisTermIn
	\def\thisterm{\thisTermIn}
\else
	\def\thisterm{Fall 2026}
\fi
%
\def\thisclass{ECE 444}

%%%%%%%%%%%%%%%%%%%%%%%%%%%%%%%%%%%%%%%%%%%%%%%%%%%%%%
%%% END SET THESE IF VARIABLES
%%%%%%%%%%%%%%%%%%%%%%%%%%%%%%%%%%%%%%%%%%%%%%%%%%%%%%
% packages and shortcuts
\input{myPackages_gr}
\usepackage[hidelinks, linkcolor=red]{hyperref}
\input{myShortcuts}
\setlength{\parindent}{0pt}
\setlength{\parskip}{6pt}

%% exam class customization
\input{exam_config}
\printselectedfalse      % full worked solutions in the key, not answers-only
\CorrectChoiceEmphasis{\color{red}}
\SolutionEmphasis{\color{red}}
\renewcommand\questionlabel{\textbf{\thequestion.}}
\renewcommand{\thepartno}{\alph{partno}}
\renewcommand{\partlabel}{(\thepartno)}

% Fillable table cell: blank in the student copy, the expected value in red on
% the KEY. Fixed-width centered columns keep the blank cells writable.
\newcolumntype{K}[1]{>{\centering\arraybackslash}p{#1}}
\newcommand{\kv}[1]{\ifkey{\color{red}#1}\fi\rule[-0.10in]{0pt}{0.34in}}

% print the answers if this is the key
\ifkey
	\printanswers % if this IS the key
\else
	\noprintanswers  %if this is NOT the key
\fi

% print the header and footer
% need to print key on top if this is the key
\ifkey
	\firstpageheader{}{\color{red}{ KEY}}{}
	\runningheader{\thisclass\hspace{1pt} -- \thisterm }{Lab \#\labNum\hspace{1pt}\color{red}{ KEY}}{\thepage\hspace{1pt} of\hspace{1pt} \numpages}
	\firstpagefooter{}{}{}
	\runningfooter{}{}{}
\else
	\firstpageheader{}{}{}
	\runningheader{\thisclass \hspace{1pt} -- \thisterm }{Lab \#\labNum\hspace{1pt}}{\thepage\hspace{1pt} of\hspace{1pt} \numpages}
	\firstpagefooter{}{}{}
	\runningfooter{}{}{}
\fi

\runningheadrule

\begin{document}
\pagestyle{headandfoot}   % exam class


%%%%%%%%%%%%%%%%%%%%%%%%%%%%%%%%%%%%%%%%%%%%%%%%%%%%%%
% title block
%%%%%%%%%%%%%%%%%%%%%%%%%%%%%%%%%%%%%%%%%%%%%%%%%%%%%%

\begin{center}
USAF ACADEMY \\
DEPARTMENT OF ELECTRICAL AND COMPUTER ENGINEERING \\
LAB \#\labNum\ (Lesson 19) - Beam Steering Lab \\
\textbf{Lab Sheet} \\
\thisterm \\[4mm]
\end{center}

\vspace{-4pt}
\textbf{Name:} \rule[-2pt]{2.6in}{0.4pt} \hfill
\textbf{Section:} \rule[-2pt]{0.9in}{0.4pt} \hfill
\textbf{Date:} \rule[-2pt]{1.1in}{0.4pt}

\vspace{4pt}

\fullwidth{\textbf{Learning Objective 3.5: } I can implement beam steering on the
ADALM-PHASER and verify the steered pattern against theory.}

\vspace{2pt}

This sheet is the turn-in document for the Lesson 19 lab. Record the source
frequency the GUI reports before you start, fill in every box, and keep the
prediction table from Lesson 18 beside you. The lab is a team assignment, so one
sheet per team.

\begin{questions}

%%%%%%%%%%%%%%%%%%%%%%%%%%%%%%%%%%%%%%%%%%%%%%%%%%%
\question \textbf{Table A -- commanded against physical angle.} Place the source
at each position on the arc, hunt for the commanded Steer Angle that maximizes
the FFT amplitude, and record it. The last column expresses the difference in
sweep grid steps, one step being the ADAR1000's $2.8125\mydeg$ phase LSB.

\begin{center}
\renewcommand{\arraystretch}{1.25}
\begin{tabular}{| l | K{1.5in} | K{1.3in} | K{1.5in} |}
\hline
\textbf{Physical angle} & \textbf{Commanded angle at peak} & \textbf{Difference} & \textbf{Difference $/\ 2.8125\mydeg$} \\\hline\hline
$0\mydeg$  & \kv{$0\mydeg \pm 2.8\mydeg$}  & \kv{$\le 2.8\mydeg$} & \kv{$\le 1$} \\\hline
$30\mydeg$ & \kv{$30\mydeg \pm 2.8\mydeg$} & \kv{$\le 2.8\mydeg$} & \kv{$\le 1$} \\\hline
$45\mydeg$ & \kv{$45\mydeg \pm 2.8\mydeg$} & \kv{$\le 2.8\mydeg$} & \kv{$\le 1$} \\\hline
\end{tabular}
\end{center}

\begin{solution}
There is no single right commanded angle here, because the physical angle set on
the arc carries its own error. The standard is one grid step: each difference
should be no larger than $2.8125\mydeg$, so the last column reads at most 1.
Grid discretization alone allows half a step, $1.41\mydeg$; protractor and
aiming error at $1\ \m$ adds about $1.5\mydeg$, HB100 drift adds a few tenths,
and room multipath moves the apparent peak by up to about a degree. Added in
quadrature those give a peak-angle uncertainty near $2.5\mydeg$. Agreement to
$0.1\mydeg$ indicates an error in how the comparison was made rather than an
unusually good measurement.
\end{solution}

%%%%%%%%%%%%%%%%%%%%%%%%%%%%%%%%%%%%%%%%%%%%%%%%%%%
\question \textbf{Table B -- phase check at $30\mydeg$.} With the peak command
applied, read the eight per-element phases out of Phase Control and compare them
against the wrapped ramp predicted in Lesson 18. At $\theta_0 = 30\mydeg$ the
progressive phase is $\Delta\phi = 176.8\mydeg \sin 30\mydeg = 88.4\mydeg$. All
values in the table are in degrees.

\begin{center}
\small
\setlength{\tabcolsep}{4pt}
\renewcommand{\arraystretch}{1.25}
\begin{tabular}{| l | K{0.55in} | K{0.55in} | K{0.55in} | K{0.55in} | K{0.55in} | K{0.55in} | K{0.55in} | K{0.55in} |}
\hline
\textbf{Element} & \textbf{1} & \textbf{2} & \textbf{3} & \textbf{4} & \textbf{5} & \textbf{6} & \textbf{7} & \textbf{8} \\\hline\hline
Predicted & \kv{$0.0$} & \kv{$88.4$} & \kv{$176.8$} & \kv{$265.2$} & \kv{$353.6$} & \kv{$82.0$} & \kv{$170.5$} & \kv{$258.9$} \\\hline
Applied & \kv{$0.00$} & \kv{$87.19$} & \kv{$177.19$} & \kv{$264.38$} & \kv{$354.38$} & \kv{$81.56$} & \kv{$171.56$} & \kv{$258.75$} \\\hline
Difference & \kv{$0.00$} & \kv{$1.21$} & \kv{$-0.39$} & \kv{$0.83$} & \kv{$-0.78$} & \kv{$0.44$} & \kv{$-1.06$} & \kv{$0.15$} \\\hline
\end{tabular}
\end{center}

\vspace{2pt}
\begin{center}
Half an LSB = \ansbox[1.3in]{$1.41\mydeg$} %
Is every difference smaller than that? \ansbox[1.0in]{yes}
\end{center}
\vspace{2pt}

\begin{solution}
The backend rounds each commanded phase to the nearest multiple of
$2.8125\mydeg$ and wraps the result into $0$ to $360\mydeg$, so no element can
miss its predicted value by more than half a step. Element 6 is the clearest
check: the unwrapped ramp value $442.0\mydeg$ has to appear as $82.0\mydeg$, and
the applied $81.56\mydeg$ is the nearest grid point to it. A difference larger
than $1.41\mydeg$ on any element means the per-element calibration offsets are
included in the reading, or the steer angle applied is not the one assumed here.
\end{solution}

%%%%%%%%%%%%%%%%%%%%%%%%%%%%%%%%%%%%%%%%%%%%%%%%%%%
\newpage
\question \textbf{Table C -- beamwidth.} Measure the half-power width of the
frozen boresight trace and of the $30\mydeg$ trace by reading the angles where
each trace falls 3 dB below its own peak, then record the peak-level difference
between the two traces.

\begin{center}
\renewcommand{\arraystretch}{1.25}
\begin{tabular}{| l | K{1.9in} | K{1.6in} |}
\hline
\textbf{Quantity} & \textbf{Measured} & \textbf{Calculated} \\\hline\hline
HPBW at $\theta_0 = 0\mydeg$ & \kv{$13.1\mydeg$, $\pm$ one grid step} & \kv{$13.0\mydeg$} \\\hline
HPBW at $\theta_0 = 30\mydeg$ & \kv{$\approx 15\mydeg$, $\pm$ one grid step} & \kv{$15.1\mydeg$} \\\hline
Peak level, $30\mydeg$ against $0\mydeg$ & \kv{$-0.5$ to $-1.5$ dB} & \kv{$-0.6$ dB} \\\hline
\end{tabular}
\end{center}

\begin{solution}
The calculated column comes from the array factor for $N = 8$ at
$d/\lambda = 0.491$. The steered beamwidth is the $1/\cos\theta_0$ broadening of
Lesson 18, $13.0\mydeg / \cos 30\mydeg = 15.0\mydeg$, and the exact array factor
gives $15.1\mydeg$. The peak drops because the array presents a projected
aperture of $Nd\cos\theta_0$ instead of $Nd$, which is
$10\mylog{\cos 30\mydeg} = -0.6$ dB for an ideal element; a real patch element
rolls off faster and usually costs a few tenths more. Measured beamwidths are
quoted to the nearest degree, because the sweep samples on a $2.8125\mydeg$ grid
and the half-power crossings land between samples.
\end{solution}

%%%%%%%%%%%%%%%%%%%%%%%%%%%%%%%%%%%%%%%%%%%%%%%%%%%
\newpage
\question \textbf{Written answers.} Three or four sentences each.

\begin{parts}

	\part Why does the commanded angle at which the trace peaks move in steps of
	about $2.8\mydeg$ rather than continuously, and what in the hardware sets that
	step size?

	\begin{solutionorlines}[2.0in]
	The ADAR1000 phase shifter is a 7-bit device, so its smallest phase increment
	is $360\mydeg / 2^7 = 2.8125\mydeg$. The sweep expresses that LSB as a
	steering resolution and visits commanded angles on the same grid, so the trace
	holds no information between grid points and the peak that is read is the grid
	point nearest the true peak. The worst case error is half a step,
	$1.41\mydeg$. No angle between two grid points can be commanded, so the peak
	moves in steps rather than continuously.
	\end{solutionorlines}

	\part Explain why a plot of received power against commanded steer angle is a
	picture of the array's pattern at all. Name the principle you are relying on,
	and say what would change if the source were moving during the sweep.

	\begin{solutionorlines}[2.3in]
	The instrument holds the source still and steps the commanded angle across its
	range, recording received power at each step. Reciprocity says the array's
	receive pattern equals its transmit pattern, so the power received from a
	fixed direction as the beam is steered traces out the same shape the array
	would radiate. The x-axis is still the commanded angle and not a measured
	arrival angle, and every point on the trace was produced by a different set of
	element phases. If the source moved during the sweep, the received power would
	change for two reasons at once, and the trace would no longer be a pattern:
	the peak would appear at whatever command happened to coincide with the
	source's closest approach to the beam.
	\end{solutionorlines}

\end{parts}

\end{questions}

\vspace{1.0em}
\noindent\textbf{Documentation:} \rule[-2pt]{0.6\linewidth}{0.4pt}

\end{document}
